
%% bare_conf.tex
%% V1.4b
%% 2015/08/26
%% by Michael Shell
%% See:
%% http://www.michaelshell.org/
%% for current contact information.
%%
%% This is a skeleton file demonstrating the use of IEEEtran.cls
%% (requires IEEEtran.cls version 1.8b or later) with an IEEE
%% conference paper.
%%
%% Support sites:
%% http://www.michaelshell.org/tex/ieeetran/
%% http://www.ctan.org/pkg/ieeetran
%% and
%% http://www.ieee.org/

%%*************************************************************************
%% Legal Notice:
%% This code is offered as-is without any warranty either expressed or
%% implied; without even the implied warranty of MERCHANTABILITY or
%% FITNESS FOR A PARTICULAR PURPOSE! 
%% User assumes all risk.
%% In no event shall the IEEE or any contributor to this code be liable for
%% any damages or losses, including, but not limited to, incidental,
%% consequential, or any other damages, resulting from the use or misuse
%% of any information contained here.
%%
%% All comments are the opinions of their respective authors and are not
%% necessarily endorsed by the IEEE.
%%
%% This work is distributed under the LaTeX Project Public License (LPPL)
%% ( http://www.latex-project.org/ ) version 1.3, and may be freely used,
%% distributed and modified. A copy of the LPPL, version 1.3, is included
%% in the base LaTeX documentation of all distributions of LaTeX released
%% 2003/12/01 or later.
%% Retain all contribution notices and credits.
%% ** Modified files should be clearly indicated as such, including  **
%% ** renaming them and changing author support contact information. **
%%*************************************************************************
\documentclass[conference]{./sty/IEEEtran}

\usepackage{cite}

% *** GRAPHICS RELATED PACKAGES ***
\usepackage[pdftex]{graphicx}
\graphicspath{{./figures/}, {./plots/}}
\usepackage[scriptsize]{subfigure}

% *** MATH PACKAGES ***
%
%\usepackage{amsmath}

% *** SPECIALIZED LIST PACKAGES ***
%
\usepackage{algorithmic}
\usepackage[usenames,dvipsnames]{xcolor}
\usepackage{ifthen}
\usepackage{amsmath}
\usepackage{amssymb}
\usepackage{acronym}
%\usepackage{url}

% *** PERSONALIZED VARS DEF ***
%
\usepackage{intero_algo_vars}

% *** Do not adjust lengths that control margins, column widths, etc. ***
% *** Do not use packages that alter fonts (such as pslatex).         ***
% There should be no need to do such things with IEEEtran.cls V1.6 and later.
% (Unless specifically asked to do so by the journal or conference you plan
% to submit to, of course. )


% correct bad hyphenation here
%\hyphenation{op-tical net-works semi-conduc-tor}

% *** COMMANDS FOR COMMENTS ***
\newboolean{showcomments}
\setboolean{showcomments}{false}
\ifthenelse{\boolean{showcomments}}
{ \newcommand{\mynote}[3]{
    \fbox{\color{#3}{#1}}
    {\small$\blacktriangleright$\texttt{\emph{\color{#3}{#2}}}$\blacktriangleleft$}}}
{ \newcommand{\mynote}[3]{}}
% One command per author:
\newcommand{\rcg}[1]{\mynote{Raziel}{#1}{Green}}
\newcommand{\er}[1]{\mynote{Etienne}{#1}{Blue}}

% *** ACRONYMS ***
\acrodef{bp}{broadcast protocol}
\acrodef{bps}{broadcast protocols}
\acrodef{Bps}{Broadcast protocols}
\acrodef{wifp}{IEEE 802.11}
\acrodef{nets}[MANETs]{mobile ad-hoc networks}


\begin{document}
%
% paper title
% Titles are generally capitalized except for words such as a, an, and, as,
% at, but, by, for, in, nor, of, on, or, the, to and up, which are usually
% not capitalized unless they are the first or last word of the title.
% Linebreaks \\ can be used within to get better formatting as desired.
% Do not put math or special symbols in the title.
\title{Towards an Adaptive Approach for Broadcast Algorithms in MANETs}


% author names and affiliations
% use a multiple column layout for up to three different
% affiliations
\author{
\IEEEauthorblockN{Raziel Carvajal-G{\'o}mez and Etienne Rivi{\`e}re}
\IEEEauthorblockA{Universit{\'e} Catholique de Louvain\\
Louvain-la-Neuve, Belgium\\
name.surname@uclouvain.be}}


% make the title area
\maketitle

\begin{abstract}
Broadcast is a fundamental operation in Mobile Ad-Hoc Networks (MANETs).
%
A large variety of broadcast protocols have been proposed to deal with network dynamics, maximize coverage and minimize the impact of collisions.
%
These protocols follow either an unstructured approach based on the control of a flooding process or a more structured approach involving the creation of a dissemination overlay.
%
The selection of the appropriate broadcast protocol for a deployed MANET is a difficult task.
%
This paper presents the ongoing work towards a novel adaptive approach to broadcast in heterogeneous MANETs.
%
This approach allows nodes to pick dynamically the most appropriate broadcast protocols for different regions of the network.
%
We present preliminary results of an interoperability mechanism to broadcast in an scenario where more than one protocol is deployed.
%
Our results show that the energy cost tends to reduce with such an approach, keeping a high coverage of nodes improving the network throughput.
\end{abstract}

% For peer review papers, you can put extra information on the cover
% page as needed:
% \ifCLASSOPTIONpeerreview
% \begin{center} \bfseries EDICS Category: 3-BBND \end{center}
% \fi
%
% For peerreview papers, this IEEEtran command inserts a page break and
% creates the second title. It will be ignored for other modes.
\IEEEpeerreviewmaketitle


\section{Introduction}
Mobile ad-hoc networks (\acs{nets}) are self-organized networks that do not relay on fixed infrastructures for their deployment.
%
Mobile nodes in such a network have wireless transceivers with a limited transmission range and forward packets to communicate with each other in a peer-to-peer fashion.
%
Application scenarios such as vehicular networks ~\cite{Slavik2014a}, monitoring of zones in danger or deploying communication infrastructures in rural areas are suitable scenarios for \acs{nets} where broadcast is a fundamental operation.
%
%\rcg{give concrete examples of systems where broadcast is important}
%

\acl{Bps} for \acs{nets} tackle the problem of selecting a set of forwarding nodes that retransmit received packets to reach every node in the network.
%
Choosing a forwarding set is not an easy task, a small set reduces the overall energy consumption but may lead in unreachable packets for other nodes.
%
On the contrarily, a bigger forwarding set increase the collisions and contentions in the network.
%
Mobility is another factor, if a forwarding set is not updated according to the velocity of nodes, obsolete choices of this set may cause lost of packets.
%
Additionally, mobility patterns may result in network partitions and heterogeneous \acs{nets} where the density of nodes varies within different sections in the area of communication.

%%
We are interested in designing adaptive approaches for \acl{bps} in \acs{nets}.
%
In our target context, adaptation refers to the ability of dynamically selecting and configuring the protocol used to broadcast.
%
Notice that the deployment of a \acl{bp} may take place either in the entire system, or more interestingly, just in a part of it.
%
To deal with the heterogeneous conditions of mobility and density in \acs{nets}, we propose an adaptive approach that converges in a system such that: on one side, overlays emerge in highly dense regions where nodes remains mostly static and on the other side, controlled flooding protocols are deployed in sparse regions where nodes move more rapidly.
%
We do not rely on a central entity with global information of the system.
%
Only with local information, nodes discover the need of adaptation and perform a collective decision towards a switch to a different protocol.
%
Our adaptive approach is provided in form of a framework with the implementation of \acl{bps} commonly found in the literature.
%
In our evaluation scenario, mobile devices are carried by humans where we use the truncated-levy walk model as a mobility pattern~\cite{Rhee:2011}.
%
The mobile device is modeled following the specifications of a transceiver, frequently used in mobile phones with \acl{wifp} capabilities.
%
Preliminary results of our simulations show that the energy cost tends to reduce with our adaptation approach, keeping a high coverage of nodes without overkilling the network bandwidth.
\section{Background}
A simple approach to broadcast is blind flooding where every node retransmit an incoming packet just once, further copies of packets previously received are ignored.
%
Blind flooding cause broadcast storms, term referring to high rate of collisions and contentions in the network~\cite{Ni:1999:BSP:313451.313525}.
%
A vast amount of protocols have been proposed to deal with broadcast storms for finding an adequate trade-off between reachability and usage of nodes resources (mainly energy consumption).
%
P Ruiz and P Bouvry ~\cite{Ruiz2015} classify \acl{bps} for \acs{nets} in two groups.
%
Protocols of the first group build and maintain an overlay by the periodic exchange of control messages that announce the existence of nodes.
%
Overlay-based approaches may form trees~\cite{796997}, clusters or connected-dominating sets (CDS), where every node is either port of a CDS or adjacent to at least one member of this set~\cite{Wu2001}.
%
In other words, CDS are forwarding sets and every other node out of a CDS act only as receiver.
%
The second group of broadcast protocols do not rely on overlays, instead, the retransmission of packets take place when predefined thresholds are reached (number of copies of received packets, timers, number of neighbors, etc) or following a stochastic approach~\cite{Kermarrec:2003}.
%
We compare our hybrid approaches to the most representative protocols of each category shown in~\cite{Ruiz2015}.
%
Blind flooding is used as lower bound of performance and as an upper bound we use a centralized implementation of the ~\cite{Esfahanian:1988:CCE:46251.46257} algorithm to find an approximation of a CDS.
%
Our work is inspired from an evaluation study of \acl{bps} for \acs{nets}~\cite{7980187} driven by a plethora of configurations of mobility and density.
%
This study reveals that overlay-based approaches perform better in a dense-mobile network where nodes remain mostly static.
%
On the contrarily, in a sparse network where nodes move faster non overlay-based approaches show better results.
%
%\rcg{Mention some adaptation approaches for broadcast algorithms. These works change the value of threshold to decide when retransmissions take place but in general, they use one protocol per deployed MANET}

\section{Interoperability in \acl{bps}}
We assume that a set $\mathbb{P}$ of broadcast protocols is in execution in a deployed mobile network of $\mathbb{N}$ nodes where each node $n_{i} \in \mathbb{N}$ runs only one protocol $p_{j} \in \mathbb{P}$ at a time.
%
Our aim is letting a node $n_{i}$ chose in an autonomous way a certain protocol $p_{j}$, this decision is based in the local information available at the node.
%
Such information may be, and it is not limited to, remaining battery, signal strength of received packets or the node relative speed.
%
But we omit the case where $n_{i}$ is running more than one protocol concurrently.

%%
When $\mathbb{P}$ contains only one protocol the reachability of messages depends on the performance of such approach.
%
Notice that in our hybrid approach we do not intent to alter the internals of any protocol in $\mathbb{P}$.
%
On the contrarily, when the cardinality of $\mathbb{P}$ is $|\mathbb{P}| > 1$ packets lost may take place.

\begin{figure}[t!]
  \centering
  \includegraphics[scale=0.20]{packets-lost-ex}
  \caption{TBD}
  \label{packets-lost-ex}
\end{figure}

%
Figure~\ref{packets-lost-ex} depicts a situation where a set of nodes $\mathbb{T}=\{t_{i} \mid i \in \{1, \ldots, 6\}\}$ run a topology-based protocol while nodes in $\mathbb{F}=\{f_{j} \mid i \in \{1, \ldots, 3\}\}$ run blind flooding.
%
Broadcast sessions initiated at $t_{1}$ form a tree such that a forwarding set composed of $t_{1}, t_{2}, t_{3}$ guarantees the message reception of every node at $\mathbb{T}$.
%
Even if nodes in $t_{4}, t_{5}, t_{6}$ are within the communication range of some nodes in $\mathbb{F}$, their role as pure receivers in the tree avoids any node in $\mathbb{F}$ hearing from messages.

\subsection{Detection of border nodes}

Algorithm~\ref{interop-alg} is a simple mechanism we propose to deal with interoperability issues between nodes running different protocols.
%%%% INTEROPERABILITY ALGORITHM %%%%

\noindent
\SetInd{0.0em}{0.60em}
%\SetAlgoNoLine
\begin{algorithm}
\footnotesize
\DontPrintSemicolon
\caption{Interoperability mechanism to deal with more than one running protocol in a MANET.}
\label{interop-alg}
\vars{
	% VARIABLES
	node\_id \textit{: string} \tcp*[1]{node identifier}
	running\_alg \textit{: string} \tcp*[1]{running algorithm}
	timeout\textit {: double} \tcp*[1]{maximum waiting time to start the procedure of border dectection}
	am\_i\_border \textit{: boolean} \tcp*[1]{says wether the local node is a border node or not}

	% SETS
	known\_algs \textit{: set} \tcp*[1]{previously known foreign protocols}

	% MAPS
	timers \textit{: dictionary} \tcp*[1]{timers of border detection procedure}
	known\_nodes \textit{: dictionary} \tcp*[1]{previously known nodes from a foreign protocol}
}

\func{ \receive{msg} \label{interop-alg:rcv-fun} }{ 
	\eIf{ msg.sender\_alg = running\_alg \Ork am\_i\_border }{
		broadcast(msg) \; \label{interop-alg:bro}
	}{
		\Switch{msg}{
			\Case{ ControlMsg }{

				known\_nodes[msg.sender\_alg] $\gets$ msg.sender\_id \; \label{interop-alg:save-sender}
				
				\If{ msg.sender\_alg $\notin$ known\_algs }{
				
					known\_algs $\gets$ known\_algs $\cup$ msg.sender\_alg \;	\label{interop-alg:save-alg}
					delay $\gets$ $new$ Random(0, timeout) \;					
					timers[msg.sender\_alg] $\gets$ $new$ Timer(delay) \; \label{interop-alg:new-timer}
					
					\TimerExpires{

						b $\gets$ $new$ BorderMsg() \; \label{interop-alg:new-border-msg}
						b.chosen\_node $\gets$ known\_nodes[msg.sender\_alg].pickOne() \label{interop-alg:new-border-msg:chosen} \;
						b.fwd\_no $\gets$ MAX\_FWD\_NO \; \label{interop-alg:new-border-msg:hops}
						b.sender\_id $\gets$ node\_id \; \label{interop-alg:new-border-msg:sender}						
						b.sender\_alg $\gets$ running\_alg \; \label{interop-alg:new-border-msg:sender-alg}
						b.foster\_by $\gets$ msg.sender\_alg \; \label{interop-alg:new-border-msg:alg}
						forward(b) \; \label{interop-alg:fwd}

					}
				}
			}
			\Case{ BorderMsg }{
			
				msg.fwd\_no $\gets$ msg.fwd\_no - 1 \; \label{interop-alg:border-msg-case:dec}
				
				\If{ msg.fwd\_no $\neq$ 0 }{
					forward(msg) \; \label{interop-alg:border-msg-case:fwd}
				}
				
				\If{ running\_alg = msg.foster\_by \And msg.foster\_by $\notin$ known\_algs }{
				
					known\_algs $\gets$ known\_algs $\cup$ msg.foster\_by \; \label{interop-alg:border-msg-case:update-known-algs-1}

					\If{ node\_id = msg.chosen\_node }{					
						am\_i\_border $\gets$ True \; \label{interop-alg:border-msg-case:select-as-border}
					}
					
				}
				
				\If{ running\_alg = msg.sender\_alg \And msg.sender\_alg $\notin$ known\_algs }{
				
					known\_algs $\gets$ known\_algs $\cup$ msg.sender\_alg \; \label{interop-alg:border-msg-case:update-known-algs-2}

					\If{ timers[msg.sender\_alg] $\neq$ $\emptyset$ }{
						cancelTimer(timers[msg.sender\_alg]) \; \label{interop-alg:border-msg-case:rm-timers}
					}
					
				}
					
			}
		}
	}
}
\end{algorithm}

%%%% ALGO %%%%
The main idea is increase the set of forwarding nodes from only one running protocol with a set we called border nodes.
%
Intuition says that the bigger the size of this set the higher the reachability of nodes among broadcast sessions.
%
Border nodes are within the transmission range of other nodes running a different protocol (this is the case of $t_{4}, t_{5}, t_{6}$ in figure~\ref{packets-lost-ex}).


%%
Our interoperability mechanism works as follows.
%
Messages in the network are triplets $msg = (\tau, \rho, \varrho)$ where $\tau$ is either a control or broadcast message, $\rho$ refers to the running protocol at the sender and $\varrho$ is the message payload.
%
Control messages are beacons, which nodes exchanged periodically, required by some protocols to operate and broadcast messages contain the information to disseminate.
%
When nodes receive a message $msg$ that corresponds to the running protocol of the sender or the current receiver is a border node (line 10), $msg$ is broadcast; notice that the function \textit{broadcast()} (line~\ref{interop-alg:bro}) hides the concrete implementation of a \acl{bp}.
%
In case the running protocols between sender and receiver differs (line 13) we distinguish between two cases:

\begin{itemize}

\item \textbf{Control message reception.}
%
This case fosters the election of border nodes.
%
We keep track of senders from a foreign algorithm with the set \textit{known\_nodes} (line~\ref{interop-alg:save-sender}), potential candidates of border nodes are within this set.
%
If receivers previously hear a beacon from a sender running a known foreign protocol, $msg$ is ignored.
%
Otherwise (line: 16), we keep a reference of the knew known protocol (line~\ref{interop-alg:save-alg}) and an event is scheduled with a timer (line~\ref{interop-alg:new-timer}).
%
Once the timer expires (line~\ref{interop-alg:timer-expires}), we create a new border message \emph{b} (line~\ref{interop-alg:new-border-msg}) with the following fields:

\begin{itemize}

\item chosen border node (line~\ref{interop-alg:new-border-msg:chosen}). For the moment this is a random choice, we can think in other methods, like piggybacking senders remaining battery to control messages for selecting a border node with the higher value of energy. 
\item number of times $b$ is forward (line~\ref{interop-alg:new-border-msg:hops})
\item sender identifier (line~\ref{interop-alg:new-border-msg:sender})
\item sender running protocol (line~\ref{interop-alg:new-border-msg:sender-alg})
\item foreign protocol that fosters the border election procedure (line~\ref{interop-alg:new-border-msg:sender-alg})

\end{itemize}
%
Now we explore the case when nodes receive border messages (line 28).

\item \textbf{Border message reception.} Receivers of this message decrease by one the number of times \textit{msg} is forwarded (line:~\ref{interop-alg:border-msg-case:dec}).
%
In case this counter has remaining hops to perform, \textit{msg} is forward (lines: 30,~\ref{interop-alg:border-msg-case:fwd}).
%
Notice that the function \textit{forward()} simply performs a retransmission within the communication range of the receiver.
%
Given that all packets are identified in a unique way further copies of previously known packets are ignored.
%
Furthermore, only one out of two possible situations take place.

%%
Either the message receiver is a potential candidate of border node (line: 32) or the receiver belongs to those nodes who foster the selection procedure (line 36).
%
When the former situation occurs, the receiver is labeled as border node only if its identifier was chosen (line~\ref{interop-alg:border-msg-case:select-as-border}).

%
The latter situation serves to cancel ongoing timers that triggered the border selection procedure (line:~\ref{interop-alg:border-msg-case:rm-timers}) by other senders, that is, nodes running a foreign algorithm. 
%

\end{itemize}


\subsection{Evaluation scenario}

We simulate a situation where a crowd of humans are carrying mobile devices with WiFi capabilities (mobile phones, tablets, etc).
%
As shown in Figure~\ref{packets-lost-ex}, the area of communication is partitioned in two regions, dense and sparse.
%
The aim of having such an arrangement is to measure the impact of discovered border nodes with our interoperability mechanism that surrounds the dense region.
%
Both regions contain 400 nodes each and approximately the same surface where the dimensions of the whole area are 120 meters of height and width.
%
Nodes are spread uniformly in both regions and their initial position is chosen in a random way.
%
The number of nodes neighbors within the sparse region varies between 1 to 10 while those nodes in the dense region may have up to 40 neighbors.
%

\smallskip
\noindent
\textbf{Broadcast protocols.}
Running protocols on each node are defined in advance and nodes keep using this protocol until the end of the simulation.
%
Nodes in the sparse area use a controlled flooding approach that do not rely on an overlay.
%
When a node receives a broadcast message for the first time it starts a timer with a waiting period chosen randomly between the interval $[0,\dots,max]$. In our experiments the threshold \emph{max} is set to 10 ms.
%
If the number of received copies per message reaches the maximum number of allowed receptions before the timer expires, the message retransmission is canceled. Otherwise, the retransmission take place.
%
This approach was first propose in~\cite{Ni:1999:BSP:313451.313525} to deal with broadcast storms.
%%
Nodes in the dense area use an approximation of a connected-dominating set called multi-point relay technique (MPR)~\cite{qayyum:inria-00072756}.

\smallskip
\noindent
\textbf{Broadcast sessions.}
A total of 200 broadcast sessions are performed with a frequency of one second.
%
The source of each sessions is chosen randomly among all nodes.
%
The content of each message is randomly generated, unique and with a length of 32~Bytes. 


\smallskip
\noindent
\textbf{Simulator.}
We deploy our simulation with OMNET++4.6~\cite{Varga:2008:OOS:1416222.1416290} and the plug-in INET3.3 to model mobile ad hoc networks.
%
This framework provides a full implementation of the IEEE 802.11g standard, we use the standard model of energy consumption.
%

\smallskip
\noindent
\textbf{Transceiver settings.}
Nodes have a transmission range of 10 m.
%
This metric depends on the wireless communication chip used.
We model power consumption using the specifications of a commercial chip found in modern mobile devices.
%
We use Broadcom BCM4329 chip~\cite{BCM4329}, which supports Wi-Fi 802.11g on the 2.4 GHz frequency band.
%
Its energy consumption per operation modes is as follows: sleep (1mW), receive (10mW) and transmit (100mW).


\smallskip
\noindent
\textbf{Mobility.}
%\subsection*{Mobility}
We use the \emph{truncated-levy walk} (LW) technique to model mobility patterns of humans carrying wireless devices~\cite{Rhee:2011}.
%
The evaluation of our interoperability mechanism is performed in a first scenario with mobility and in a second scenario where nodes remain static.
%
In the scenario with mobility, nodes at the dense region move at a speed of 0.5 m/s and the speed of those at the sparse region is 2.0 m/s; these values are based on an estimation of the preferred walking speed of humans in~\cite{Browning390}.
%
The LW model generates traces of frequent short walks and occasional rides to distant locations.
%
Nodes move within their region of density in the communication area without getting out of their designated region.

	
\subsection{Preliminary Results}
\rcg{TBD...}

\begin{figure}[!t]
	\begin{center}
		\subfigure[Overall energy consumption]{
			\includegraphics[width=\linewidth]{energy_consump_static}
			\label{plots:energy-consump}
		}
		\subfigure[Relative error of received messages]{
			\includegraphics[width=\linewidth]{relative_error_rcv_msgs_static}
			\label{plots:re-rcv-msgs}
		}
				\subfigure[Network coverage]{
			\includegraphics[width=\linewidth]{netowrk_coverage_static}
			\label{plots:net-coverage}
		}
		\caption{Preliminary results of broadcast metrics in the static scenario.} 
		\label{fig:outbw}
	\end{center}
\end{figure}

\begin{figure}[!t]
  \centering
  \includegraphics[trim = 0 0 0 9mm, clip, scale=0.4]{fwd_types_dense_area_static}
  \caption{Forwarding sets grouped by type of retransmission}
  \label{plots:fwd-sets-type}
\end{figure}
\section{Future work and Conclusion}
Selecting what is the appropriate broadcast algorithm for a given MANET deployment scenario is a daunting task.
%
The performance of a given approach depends on its deployment context.
This makes hard to find the right balance between reliability and efficiency.
%
We present preliminary results of an hybrid approach for broadcast protocols in MANETs where an overlay is deployed in dense regions of the network and controlled flooding techniques run in sparse regions.

%
As future work, we will focus in the specifications of adaptation rules for lettings nodes decide which broadcast protocol fits better in the network.
%
This decision may be collective or individual but completely decentralized.
%
We rely on the local meta data available in nodes such as: the signal strength of received packets, remaining energy or rate of received packets.
%
Mainly, nodes velocity and density are two transversal aspects to take into consideration.
%


% use section* for acknowledgment
\section*{Acknowledgment}
The research leading to these results has received funding from CHIST-ERA under project DIONASYS, from the Swiss National Science Foundation (SNSF) under grant 155249, and from the French National Research Agency (ANR) under grant ANR-14-CHR2-0004.

% references section
\bibliographystyle{./sty/IEEEtran}
\bibliography{main}


% that's all folks
\end{document}


